% Computer Science A --- Web fundamentals and Gradle (practice bank).
% Formatted from raw source in this file.

\runningheader{Computer Science A}{Web \& Gradle practice}{Page \thepage\ of \numpages}

\begingradingrange{csawebtest}

\uplevel{\par\textbf{Part 1: HTML}\par}

\question[4]
Which HTML version introduced semantic elements?
\begin{choices}
  \choice HTML 4
  \choice XHTML
  \CorrectChoice HTML5
  \choice XML
\end{choices}
\begin{solution}
  \textbf{HTML5} (2014) added semantic landmarks such as \lstinline[style=htmlinline]|<nav>|, \lstinline[style=htmlinline]|<article>|, \lstinline[style=htmlinline]|<section>|, and \lstinline[style=htmlinline]|<header>|. HTML 4 and XHTML relied heavily on generic \lstinline[style=htmlinline]|<div>| wrappers; XML is a markup meta-language, not an HTML version.
\begin{lstlisting}[style=html]
<header><h1>Course site</h1></header>
<nav aria-label="Main"><a href="/">Home</a></nav>
<main><article><h2>Week 3</h2></article></main>
\end{lstlisting}
\end{solution}

\question[4]
What is the basic structure of an HTML document?
\begin{choices}
  \choice \texttt{head}, \texttt{body}, \texttt{footer}
  \CorrectChoice \texttt{html}, \texttt{head}, \texttt{body}
  \choice \texttt{body}, \texttt{div}, \texttt{span}
  \choice \texttt{meta}, \texttt{link}, \texttt{title}
\end{choices}
\begin{solution}
  Every HTML document is wrapped in \lstinline[style=htmlinline]|<html>|. \lstinline[style=htmlinline]|<head>| holds metadata; \lstinline[style=htmlinline]|<body>| holds visible content. \texttt{footer} is an optional semantic element inside \texttt{body}, not a required top-level sibling of \texttt{head}.
\begin{lstlisting}[style=htmlinline]
<!DOCTYPE html>
<html lang="en">
  <head><title>Lab</title></head>
  <body><p>Visible content</p></body>
</html>
\end{lstlisting}
\end{solution}

\question[4]
Which tag creates a hyperlink?
\begin{choices}
  \choice \texttt{<img>}
  \choice \texttt{<link>}
  \CorrectChoice \texttt{<a>}
  \choice \texttt{href}
\end{choices}
\begin{solution}
  The anchor element \lstinline[style=htmlinline]|<a>| defines a hyperlink. \lstinline[style=htmlinline]|href| is an \emph{attribute} on \lstinline[style=htmlinline]|<a>|, not a tag. \lstinline[style=htmlinline]|<link>| connects external resources in \lstinline[style=htmlinline]|<head>| (stylesheets, icons), not inline navigation links.
\begin{lstlisting}[style=htmlinline]
<a href="/docs">Read the documentation</a>
\end{lstlisting}
\end{solution}



\newpage




\question[4]
Which element is block-level?
\begin{choices}
  \choice \texttt{<span>}
  \choice \texttt{<img>}
  \CorrectChoice \texttt{<div>}
  \choice \texttt{<strong>}
\end{choices}
\begin{solution}
  \lstinline[style=htmlinline]|<div>| is a generic block container: it starts on a new line and spans the full width of its parent. \lstinline[style=htmlinline]|<span>| and \lstinline[style=htmlinline]|<strong>| are inline; \lstinline[style=htmlinline]|<img>| is a replaced inline element.
\begin{lstlisting}[style=htmlinline]
<div>Block (new line, full width)</div>
<span>Inline</span><strong>inline</strong>
\end{lstlisting}
\end{solution}

\question[4]
What is the purpose of semantic elements?
\begin{choices}
  \choice Improve styling
  \CorrectChoice Improve code readability and SEO
  \choice Reduce file size
  \choice Replace CSS
\end{choices}
\begin{solution}
  Semantic tags describe \emph{meaning} (\emph{this is navigation}, \emph{this is an article}), helping developers, search crawlers, and assistive technologies understand page structure. Styling still requires CSS; file size is largely unchanged.
\begin{lstlisting}[style=html]
<article>
  <h2>Release notes</h2>
  <p>Search engines and screen readers identify this as main content.</p>
</article>
\end{lstlisting}
\end{solution}

\question[4]
Which HTML element is used for navigation?
\begin{choices}
  \choice \texttt{<section>}
  \choice \texttt{<article>}
  \CorrectChoice \texttt{<nav>}
  \choice \texttt{<header>}
\end{choices}
\begin{solution}
  \lstinline[style=htmlinline]|<nav>| is the dedicated landmark for major navigation blocks. \lstinline[style=htmlinline]|<header>| is a page or section header; \lstinline[style=htmlinline]|<section>| and \lstinline[style=htmlinline]|<article>| group content but do not imply navigation.
\begin{lstlisting}[style=htmlinline]
<nav aria-label="Primary">
  <a href="/">Home</a>
  <a href="/labs">Labs</a>
</nav>
\end{lstlisting}
\end{solution}



\newpage



\question[4]
Which attribute enables form validation?
\begin{choices}
  \choice \texttt{action}
  \choice \texttt{method}
  \CorrectChoice \texttt{required}
  \choice \texttt{value}
\end{choices}
\begin{solution}
  The boolean \lstinline[style=htmlinline]|required| attribute tells the browser the field must be filled before submit. \lstinline[style=htmlinline]|action| and \lstinline[style=htmlinline]|method| control where and how the form is sent; \lstinline[style=htmlinline]|value| sets the default field content.
\begin{lstlisting}[style=htmlinline]
<input type="email" id="email" name="email" required>
\end{lstlisting}
\end{solution}

\question[4]
What does \lstinline[style=jsinline]|localStorage| provide?
\begin{choices}
  \choice Server-side storage
  \choice Temporary memory (cleared when the tab closes)
  \CorrectChoice Persistent client-side storage
  \choice Cache only (never writable by scripts)
\end{choices}
\begin{solution}
  \lstinline[style=jsinline]|localStorage| stores key--value strings in the browser for the origin until explicitly cleared. Data survives tab closes and browser restarts (unlike \lstinline[style=jsinline]|sessionStorage|). It is not server-side and is writable from JavaScript.
\begin{lstlisting}[style=js]
localStorage.setItem('theme', 'dark');
console.log(localStorage.getItem('theme')); // "dark"
\end{lstlisting}
\end{solution}

\newpage

\uplevel{\par\textbf{Part 2: CSS}\par}

\question[4]
Which language is used for styling web pages?
\begin{choices}
  \choice HTML
  \choice JavaScript
  \CorrectChoice CSS
  \choice SQL
\end{choices}
\begin{solution}
  \textbf{Cascading Style Sheets (CSS)} control layout, color, typography, and responsive behavior. HTML defines structure; JavaScript adds behavior; SQL queries databases.
\begin{lstlisting}[style=css]
h1 { color: navy; font-size: 2rem; }
.card { padding: 1rem; border-radius: 8px; }
\end{lstlisting}
\end{solution}

\question[4]
Which selector targets an element by ID?
\begin{choices}
  \choice \texttt{.class}
  \CorrectChoice \texttt{\#id}
  \choice \texttt{element}
  \choice \texttt{*}
\end{choices}
\begin{solution}
  The ID selector uses a hash prefix: \lstinline[style=cssinline]|#save-btn|. Class selectors use a dot; type selectors use the tag name; \lstinline[style=cssinline]|*| matches every element.
\begin{lstlisting}[style=htmlinline]
<button id="save-btn">Save</button>
\end{lstlisting}
\begin{lstlisting}[style=css]
#save-btn { background: green; color: white; }
\end{lstlisting}
\end{solution}

\question[4]
What does the CSS box model consist of?
\begin{choices}
  \CorrectChoice Margin, border, padding, and content
  \choice Width and height only
  \choice Color and font properties
  \choice Grid and flex tracks
\end{choices}
\begin{solution}
  From outside in: \textbf{margin} (outer spacing), \textbf{border}, \textbf{padding} (inner spacing), then \textbf{content}. \lstinline[style=cssinline]|box-sizing: border-box| includes padding and border inside the declared width.
\begin{lstlisting}[style=css]
.card {
  width: 200px;
  padding: 16px;
  border: 2px solid gray;
  margin: 12px;
}
\end{lstlisting}
\end{solution}



\newpage



\question[4]
Which layout system is one-dimensional?
\begin{choices}
  \choice Grid
  \choice Float
  \CorrectChoice Flexbox
  \choice Table
\end{choices}
\begin{solution}
  \textbf{Flexbox} lays out items along a single main axis (row \emph{or} column). \textbf{CSS Grid} is two-dimensional (rows and columns together). Float and table layouts are older patterns; neither is the modern one-axis tool.
\begin{lstlisting}[style=css]
.toolbar {
  display: flex;
  justify-content: space-between;
  align-items: center;
}
\end{lstlisting}
\end{solution}

\question[4]
Which CSS feature enables responsive design?
\begin{choices}
  \choice Flexbox alone
  \CorrectChoice Media queries
  \choice Animations
  \choice Pseudo-elements
\end{choices}
\begin{solution}
  \textbf{Media queries} apply styles when viewport width, orientation, or user preferences match a condition. Flexbox helps layout but does not by itself switch rules at breakpoints.
\begin{lstlisting}[style=css]
.container { width: 100%; }
@media (min-width: 768px) {
  .container { width: 90%; max-width: 1200px; }
}
\end{lstlisting}
\end{solution}

\newpage

\uplevel{\par\textbf{Part 3: JavaScript}\par}

\question[4]
What is JavaScript primarily used for?
\begin{choices}
  \choice Styling pages
  \choice Database management
  \CorrectChoice Client-side interactivity
  \choice Server hosting
\end{choices}
\begin{solution}
  JavaScript runs in the browser (and on servers via Node.js) to respond to events, update the DOM, and call APIs. CSS handles styling; databases and hosting are separate infrastructure concerns.
\begin{lstlisting}[style=js]
document.getElementById('btn').addEventListener('click', () => {
  alert('Button clicked');
});
\end{lstlisting}
\end{solution}

\question[4]
Which keyword declares a block-scoped variable?
\begin{choices}
  \choice \lstinline[style=jsinline]|var|
  \CorrectChoice \lstinline[style=jsinline]|let|
  \choice \lstinline[style=jsinline]|static|
  \choice \lstinline[style=jsinline]|define|
\end{choices}
\begin{solution}
  \lstinline[style=jsinline]|let| and \lstinline[style=jsinline]|const| are block-scoped. \lstinline[style=jsinline]|var| is function-scoped and hoisted. \lstinline[style=jsinline]|static| belongs to class syntax in ES6 classes; \lstinline[style=jsinline]|define| is not a JavaScript keyword.
\begin{lstlisting}[style=js]
if (true) {
  let count = 0;
  var legacy = 1;
}
// count is not in scope here; legacy leaks to outer function/global
\end{lstlisting}
\end{solution}

\question[4]
Which array method transforms elements?
\begin{choices}
  \choice \lstinline[style=jsinline]|.filter()|
  \choice \lstinline[style=jsinline]|.reduce()|
  \CorrectChoice \lstinline[style=jsinline]|.map()|
  \choice \lstinline[style=jsinline]|.forEach()|
\end{choices}
\begin{solution}
  \lstinline[style=jsinline]|.map()| applies a function to each element and returns a \emph{new} array of transformed values. \lstinline[style=jsinline]|.filter()| selects a subset; \lstinline[style=jsinline]|.reduce()| accumulates to one value; \lstinline[style=jsinline]|.forEach()| runs side effects without returning a mapped array.
\begin{lstlisting}[style=js]
const nums = [1, 2, 3];
const doubled = nums.map(n => n * 2); // [2, 4, 6]
\end{lstlisting}
\end{solution}



\newpage



\question[4]
What does \lstinline[style=jsinline]|this| refer to in JavaScript?
\begin{choices}
  \choice The global scope only
  \CorrectChoice The current execution context
  \choice A reserved variable name for loops
  \choice Always the function that declared it
\end{choices}
\begin{solution}
  \lstinline[style=jsinline]|this| is bound by \emph{how} a function is called: as a method it often refers to the object before the dot; in strict mode at top level it may be \lstinline[style=jsinline]|undefined|; arrow functions inherit \lstinline[style=jsinline]|this| from the enclosing scope.
\begin{lstlisting}[style=js]
const user = {
  name: 'Ada',
  greet() { console.log(this.name); }
};
user.greet(); // "Ada" - this is user
\end{lstlisting}
\end{solution}

\question[4]
Which feature handles asynchronous operations?
\begin{choices}
  \choice \lstinline[style=jsinline]|for| loop
  \choice Callbacks
  \CorrectChoice Promises
  \choice \lstinline[style=jsinline]|async|/\lstinline[style=jsinline]|await|
\end{choices}
\begin{solution}
  A \textbf{Promise} is the standard object representing deferred work (\emph{pending}, \emph{fulfilled}, \emph{rejected}). Callbacks and \lstinline[style=jsinline]|async|/\lstinline[style=jsinline]|await| are related patterns---\lstinline[style=jsinline]|async|/\lstinline[style=jsinline]|await| is syntactic sugar built on Promises---but Promises are the core reusable abstraction in modern JavaScript.
\begin{lstlisting}[style=js]
fetch('/api/user')
  .then(res => res.json())
  .then(data => console.log(data))
  .catch(err => console.error(err));
\end{lstlisting}
\end{solution}

\newpage

\uplevel{\par\textbf{Part 4: Gradle}\par}

\question[4]
Which repositories can Gradle use for resolving dependencies? \emph{Select all that apply.}
\begin{checkboxes}
  \CorrectChoice Maven Central repository
  \choice Docker Hub
  \CorrectChoice Ivy repositories
  \CorrectChoice Local repositories
\end{checkboxes}
\begin{solution}
  Gradle resolves artifacts from Maven-compatible repos (including Maven Central), Ivy repos, and the local cache (typically \lstinline[style=bashinline]|~/.gradle/caches| or a declared local directory). Docker Hub distributes container images, not Java library coordinates for Gradle dependency resolution.
\begin{lstlisting}[style=java]
repositories {
    mavenCentral()
    ivy { url "https://example.com/ivy" }
    mavenLocal()
}
\end{lstlisting}
\end{solution}

\question[4]
For Gradle projects, a unit test library like JUnit would go in which section?
\begin{choices}
  \choice \texttt{plugins}
  \choice \texttt{repositories}
  \CorrectChoice \texttt{dependencies}
  \choice \texttt{tasks}
\end{choices}
\begin{solution}
  Library coordinates (for example \lstinline[style=javainline]|testImplementation 'org.junit.jupiter:junit-jupiter:5.10.0'|) belong in the \textbf{dependencies} block. \texttt{plugins} apply Gradle plugins; \texttt{repositories} declare where to download artifacts; \texttt{tasks} define or configure build steps.
\begin{lstlisting}[style=java]
dependencies {
    testImplementation 'org.junit.jupiter:junit-jupiter:5.10.0'
}
\end{lstlisting}
\end{solution}

\question[4]
What are Gradle build scripts?
\begin{choices}
  \CorrectChoice Script files that define Gradle \textbf{projects} and \textbf{tasks}; a project may represent a library JAR, web application, or assembled deliverable
  \choice Script files that list dependencies stored as JARs in the project root package only
  \choice XML project descriptors with default directories such as \texttt{src/main/java} (Maven \texttt{pom.xml} style)
  \choice JVM instruction files compiled with \lstinline[style=javainline]|javac| into \texttt{.class} files
\end{choices}
\begin{solution}
  Gradle builds are driven by scripts (commonly \lstinline[style=bashinline]|build.gradle| or \lstinline[style=bashinline]|build.gradle.kts|). Each build has one or more \textbf{projects}; each project registers \textbf{tasks} (compile, test, jar, etc.). Dependencies are declared in the script but downloaded to Gradle's cache, not checked into the repo root. The XML description matches Maven, not Gradle; build scripts are not Java source compiled by \lstinline[style=javainline]|javac|.
\end{solution}



\newpage




\question[4]
What is Gradle?
\begin{choices}
  \choice A dependency management tool only
  \choice The Java compiler
  \choice A container orchestrator
  \CorrectChoice A build automation tool
\end{choices}
\begin{solution}
  \textbf{Gradle} automates compiling, testing, packaging, and publishing software. It \emph{includes} dependency management but also orchestrates tasks and plugins across projects. The Java compiler is part of the JDK; container orchestration (for example Kubernetes) is a separate concern.
\begin{lstlisting}[style=bash]
./gradlew build    # compile, test, and assemble artifacts
./gradlew test     # run unit tests
\end{lstlisting}
\end{solution}

\endgradingrange{csawebtest}
